\section[Betweenness Ranking Approximation]{Betweenness Ranking Approximation}\label{sec:betweenness-ranking-approximation}

The betweenness centrality of a node measures its importance as a bridge or bottleneck based on the fraction of all shortest paths passing through
 it~\parencite{freeman1977set,brandes_faster_2001}.
 Given an unweighted graph $G=(V,E)$, let $\sigma_{s,t}$ denote the total number of shortest paths from node $s$ to node $t$,
  and let $\sigma_{s,t}(v)$ be the number of those paths that pass through $v \in V$.
The exact betweenness centrality $C_B(v)$ of a vertex $v$ is formally defined as:

\begin{equation}
    C_B(v) = \sum_{s \neq v \neq t \in V} \frac{\sigma_{s,t}(v)}{\sigma_{s,t}} \label{eq:betweenness-def}
\end{equation}

where $\frac{\sigma_{s,t}(v)}{\sigma_{s,t}} = 0$ if $\sigma_{s,t} = 0$~\parencite{brandes2001faster}. In literature and practical network libraries, normalized variants are commonly used to ensure scale-invariance across different graph sizes.

The standard normalization factor depends on whether the source and target node pairs can include $v$ in their normalization pool. For instance, the Julia \texttt{Graphs.jl} library normalizes by dividing by the total number of pairs excluding $v$:

\begin{equation}
    C_B(v) = \frac{1}{(n-1)(n-2)}\sum_{s \neq v \neq t \in V} \frac{\sigma_{s,t}(v)}{\sigma_{s,t}} \label{eq:betweenness-julia}
\end{equation}

Alternatively, the adaptive approximation algorithm \textsc{Kadabra} treats betweenness centrality probabilistically~\parencite{borassi_kadabra_2019}.
It frames the metric as the exact probability that a shortest path $\pi$ between two uniformly selected random nodes $s$ and $t$
 passes through $v$~\parencite{borassi2019kadabra}.
Following the convention where the sum includes cases where $\sigma_{s,t}(v)=0$ whenever $s=v$ or $t=v$,
the normalized centrality is expressed as:

\begin{equation}
    C_B(v) = \frac{1}{n(n-1)} \sum_{s \neq t \in V} \frac{\sigma_{s,t}(v)}{\sigma_{s,t}} \label{eq:betweenness-kadabra}
\end{equation}


\subsection{Betweenness Ranking Approximation Problem}
 Given the prohibitive $O(|V||E|)$ exact computation cost of Brandes' algorithm on massive networks,
 the goal of the \emph{Betweenness Ranking Approximation Problem} is to learn a size-invariant function $f: V \rightarrow \mathbb{R}$
that maps each node to a real-valued score $f(v)$ such that the induced relative ordering maximizes the ordinal correlation 
with the true centrality hierarchy:

\begin{equation}
    \max_{f} \tau \left( \mathcal{R}(f(V)), \mathcal{R}(C_B(V)) \right) \label{eq:ranking-objective}
\end{equation}

Where $\mathcal{R}(\cdot)$ represents the rank operator, and $\tau$ denotes a ranking index such as the Kendall Tau correlation coefficient.

This formulation allows the estimation to be generalized natively across edge-weighted or directed graph structures.
 For directed representations, the flow maps independent path dynamics from incoming and outgoing topologies before scoring fusion ($w_{u,v}=1$ for structural baselines).

For later structural representation, we define the $m$-th order multi-hop degree mass of a node vector as:
\begin{equation}
    d^{(m)}=\left(\sum_{k=0}^{m}A^{k}\right)d
\end{equation}
as the cumulative structural neighborhood volume within $m$ hops, where $A$ represents the graph adjacency matrix and $d$ is the standard degree vector. In strictly local configurations, $m=1$ holds.

The problem is equivalent to preserving topological dominance hierarchies without minimizing absolute regression errors.
 For high-diameter configurations like spatial road networks, structural regularities produce strict discrete symmetry boundaries,
  causing standard scale-free embeddings to experience over-smoothing.
   By shifting the inputs to size-invariant degree masses, long-range structural bottlenecks are captured efficiently
    within a lightweight parameter footprint.

\begin{align}
    \mathcal{L}(s_u, s_v, y) & = \max(0, 1 - y \cdot (s_u - s_v))
\end{align}

Since exact deep architectures for this structural proxy require representative generative training,
 formulations leverage synthetic hyperplanes such as the hyperbolic random graph model to calibrate power-law
  degree variations outside standard scale-free assumptions prior to message passing.